% 2de graad, reeks Ruimtemeetkunde.
% Saenredam: een witte kerk vol evenwijdige en snijdende lijnen en vlakken.
% Motief: een kubus met twee kruisende ribben: niet evenwijdig en zonder snijpunt.
\def\omslagbestand{omslag-assendelft.jpg}
\def\omslagbron{Pieter Saenredam, \textit{Interieur van de Sint-Odulphuskerk in Assendelft}, 1649 (publiek domein)}
\def\omslagrand{10.6}
\def\omslagmotief{%
  \randlijn
  \begin{scope}[shift={(1,6.4)}]
  \draw[omslaglijn=.7pt,dashed] (16,5.5) -- (18.4,5.5) (16,5.5) -- (16,7.9) (16,5.5) -- (15,4.8);
  \draw[omslaglijn=1.1pt,line join=round] (15,4.8) rectangle (17.4,7.2);
  \draw[omslaglijn=1.1pt,line join=round] (15,7.2) -- (16,7.9) -- (18.4,7.9) -- (18.4,5.5) -- (17.4,4.8)
    (17.4,7.2) -- (18.4,7.9);
  \draw[fel,line width=2.2pt,preaction={draw=accent,line width=3.6pt}] (14.4,4.8) -- (18,4.8);
  \draw[fel,line width=2.2pt,preaction={draw=accent,line width=3.6pt}] (18.4,5) -- (18.4,8.6);
  \node[omslaglabel,anchor=south] at (18.4,8.7) {kruisend};
  \end{scope}
}
